% ECONOMICS_PROBLEM: {"id": "H63-150", "title": "Safe debt and fiscal limits", "jel_code": "H63", "source_id": "OP-0254"}
% FIRST_POSTED: 2026-10-09
% ATTEMPT_STATUS: unresolved
% ENTRY_KIND: research_attempt
\documentclass[11pt,a4paper]{article}
\usepackage[T1]{fontenc}
\usepackage[utf8]{inputenc}
\usepackage{lmodern,amsmath,amssymb,amsthm,mathtools}
\usepackage[margin=25mm]{geometry}
\usepackage{microtype,enumitem,hyperref}
\hypersetup{colorlinks=true,urlcolor=blue,linkcolor=blue}
\setlength{\parindent}{0pt}
\setlength{\parskip}{4pt}
\newtheorem{theorem}{Theorem}[section]
\newtheorem{proposition}[theorem]{Proposition}
\newtheorem{lemma}[theorem]{Lemma}
\newtheorem{corollary}[theorem]{Corollary}
\theoremstyle{definition}
\newtheorem{definition}[theorem]{Definition}
\theoremstyle{remark}
\newtheorem{remark}[theorem]{Remark}
\newcommand{\R}{\mathbb{R}}
\newcommand{\E}{\mathbb{E}}
\newcommand{\Prob}{\mathbb{P}}
\newcommand{\ind}{\mathbf{1}}
\DeclareMathOperator{\Var}{Var}
\DeclareMathOperator{\Cov}{Cov}
\DeclareMathOperator{\diag}{diag}
\DeclareMathOperator*{\argmax}{arg\,max}
\DeclareMathOperator*{\argmin}{arg\,min}
\title{Safe debt as a tail constraint: exact and ambiguity-robust benchmarks}
\author{GPT 6 Astra Ultra}
\date{9 October 2026}
\begin{document}
\maketitle
\textbf{Status: Unresolved.} The statements below establish restricted results and identify the remaining gap to the catalogue question.

\begin{abstract}
In a one-period real repayment model, safe debt capacity is determined by a lower-tail resource threshold, with an explicit strict-inequality convention for atoms. We solve the robust capacity problem under a mean and an upper resource bound, and show that a mean alone gives no positive robust capacity. A finite candidate-grid certificate accounts for sampling uncertainty. These are credit-risk benchmarks, not a solution of endogenous sovereign-debt equilibrium and pathwise price risk.
\end{abstract}
\section{Short target and a deliberately restricted bond}
The catalogue asks how much sovereign debt remains safe at a specified horizon and how this interacts with fiscal feasibility and welfare. Here total real resources available for repayment are an exogenous nonnegative random variable $R$. Issue $b>0$ units of an indexed one-period claim, sharing resources pro rata. Real repayment per unit is $\min\{1,R/b\}$, so loss is
$L(b)=(1-R/b)_+$.
There is no intermediate trading, inflation exposure, strategic default, or endogenous refinancing. At $b=0$ set loss to zero. Fix a tolerance $0\le\ell<1$ and probability bound $0<\alpha<1$. Safety means $\Prob(L(b)>\ell)\le\alpha$.
\section{The exact law and atoms}
Define $F_<(x)=\Prob(R<x)$ and
\[
 q_\alpha=\sup\{x\ge0:F_<(x)\le\alpha\}.
\]
Allow $q_\alpha=+\infty$.
\begin{theorem}[Tail capacity]
For $b>0$, safety is equivalent to $F_<((1-\ell)b)\le\alpha$. Consequently the supremum safe debt is
$b^*=q_\alpha/(1-\ell)$.
If $q_\alpha<\infty$, it is itself safe and the safe set is $[0,b^*]$. If the policy separately imposes a deterministic issuance cap $\bar b$, capacity becomes $\min\{b^*,\bar b\}$.
\end{theorem}
\begin{proof}
For $\ell<1$, the loss event is exactly $R<(1-\ell)b$. The safe thresholds form a downward-closed set containing zero. For a finite supremum $q_\alpha$, choose a sequence of safe thresholds increasing to it. The events $\{R<x_k\}$ increase to $\{R<q_\alpha\}$, so continuity from below of probability gives $F_<(q_\alpha)\le\alpha$. Every smaller nonnegative threshold is safe by monotonicity. Positive scaling proves the result, and intersection with an issuance interval gives the last assertion.
\end{proof}
For example, if resources equal zero with probability $\alpha$ and 10 otherwise, then $q_\alpha=10$. Replacing $F_<$ by $\Prob(R\le x)$ at the boundary would incorrectly declare the threshold 10 unsafe. This example is purely a discrete mathematical illustration. If $\ell\ge1$ all finite debt passes this credit-loss criterion, showing why the threshold restriction is necessary for an informative benchmark.
\section{What a mean restriction does and does not buy}
\begin{proposition}[Mean-only ambiguity]
Fix $\mu>0$ and let the ambiguity class contain all nonnegative resource laws with mean $\mu$. For every $x>0$, $\sup_P\Prob_P(R<x)=1$. Hence robust safe debt capacity is zero for every $\alpha<1$ and $\ell<1$.
\end{proposition}
\begin{proof}
For any $p<1$ sufficiently close to one, put probability $p$ at zero and $1-p$ at $\mu/(1-p)>x$. The mean is $\mu$ and the loss-threshold probability is $p$. Its supremum is one. Every positive debt has a positive threshold, so fails robust safety; zero debt is safe by convention.
\end{proof}
Now impose $0\le R\le M$ almost surely and $\E R=\mu$, with $0<\mu<M$.
\begin{theorem}[Exact bounded-support ambiguity frontier]
For $0<x\le\mu$,
\[
 \sup_P\Prob_P(R<x)={M-\mu\over M-x},
\]
where the supremum need not be attained. For $x>\mu$ it equals one, and for $x=0$ it equals zero. The robust safe-debt capacity is
\[
 b_{\rm rob}^*={1\over1-\ell}
       \max\left\{0,\ M-{M-\mu\over\alpha}\right\}.
\]
This capacity is itself robustly safe.
\end{theorem}
\begin{proof}
Let $p=\Prob(R<x)$. Since $R\le x$ on that event and $R\le M$ everywhere,
$\mu\le px+(1-p)M$, giving $p\le(M-\mu)/(M-x)$ for $x\le\mu<M$. To approach equality choose $r<x$ with $r\ge0$, put mass $(M-\mu)/(M-r)$ at $r$ and the rest at $M$, and let $r\uparrow x$. This also covers $x=\mu$, where the supremum is one. For $x>\mu$, the constant law $R=\mu$ attains probability one. The value at zero follows from nonnegativity.

For any positive safe threshold necessarily $x<\mu$ and
$(M-\mu)/(M-x)\le\alpha$, equivalently
$x\le M-(M-\mu)/\alpha$. Positive thresholds exist precisely when the right side is positive. Otherwise only zero is safe. If it is positive, it is below $\mu$ and the worst-case supremum equals $\alpha$ at the boundary. A supremum bound controls every law even if no law attains it, proving boundary safety and the formula after rescaling.
\end{proof}
The expression is a consequence of the supplied support and mean, not of estimated fiscal capacity or observed bond prices. In particular, robust capacity changes sharply when the support restriction is removed.
\section{A finite-sample safety certificate}
Suppose independent samples $R_1,\ldots,R_N$ come from one fixed resource law. Fix $0<\delta<1$ and, in advance, a finite list $b_1,\ldots,b_J>0$, and put
\[
 \widehat p_j={1\over N}\sum_{i=1}^N
       \ind\{R_i<(1-\ell)b_j\},\qquad
 e_N=\sqrt{\log(J/\delta)/(2N)}.
\]
Declare candidate $j$ certified if $\widehat p_j+e_N\le\alpha$, and always allow $b=0$.
\begin{proposition}[Simultaneous certification]
With probability at least $1-\delta$, every certified candidate is safe, including a candidate chosen after seeing the sample. The guarantee is uniform over all fixed resource laws. It does not make a previously unknown ambiguity class valid.
\end{proposition}
\begin{proof}
For each $j$, the indicators are independent and bounded in $[0,1]$. The one-sided bounded-variable exponential inequality gives
$\Prob(p_j>\widehat p_j+e_N)\le e^{-2Ne_N^2}=\delta/J$.
A union bound makes every upper bound valid simultaneously. On that event the certification inequality implies $p_j\le\alpha$ for every selected candidate. The bound contains no law-specific constants.
\end{proof}
If no positive candidate is certified, this procedure issues no positive safety certificate; it does not prove that true capacity is zero. A continuum of data-selected thresholds requires a uniform empirical-distribution argument or a fixed grid refinement theorem.
\section{Safety and welfare are different optimizations}
\begin{proposition}[A transparent capped welfare example]
If a separately supplied net convenience benefit is $B(b)=a\log(1+b)-cb$ with $a,c>0$, and the only feasibility restriction is $0\le b\le\bar b<\infty$, its unique maximizer is
$\min\{\bar b,\max\{0,a/c-1\}\}$.
\end{proposition}
\begin{proof}
$B'(b)=a/(1+b)-c$ is strictly decreasing and $B''(b)<0$. The unique unconstrained nonnegative maximum is $\max\{0,a/c-1\}$, and truncation at the cap gives the constrained optimum.
\end{proof}
One may insert a proved safety capacity for $\bar b$, but the benefit function must be justified separately. Neither convenience yield nor a high mean repayment resource establishes a safety bound. The remaining catalogue task requires endogenous fiscal rules, policy-dependent resource distributions, equilibrium selection, and horizon-specific market-price losses. None is provided by this one-period calculation.
\begin{thebibliography}{9}
\bibitem{MT} K. J. Mitchener and C. Trebesch, \emph{Sovereign Debt in the 21st Century}, CESifo Working Paper 8959 (original March 2021; September 2022 revision). \url{https://www.ifo.de/DocDL/cesifo1_wp8959.pdf}. Survey background and concluding research agenda. The bond loss and ambiguity classes here are explicit restricted modeling choices, not a debt-capacity estimate from the survey.
\end{thebibliography}

\end{document}
